\section{方法}
\label{sec:method}

\subsection{问题建模}

给定一个二维部署区域R和J个雷达节点，需要优化每个雷达的位置。令X = {x\_1, x\_2, ..., x\_J}表示雷达的物理位置，其中x\_j = (x\_j, y\_j)。

定义任务点集合T = {t\_1, t\_2, ..., t\_M}，其中t\_m = (x\_m, y\_m)。每个任务点有权重w\_m，默认值为1.0。

\textbf{目标函数1 - 感知效能：}
\begin{equation}
    \text{ECR} = \frac{1}{W} \sum_{m=1}^{M} w_m \cdot I\left(P_{\text{joint}}(m) \geq P_{\text{th}}\right)
\end{equation}
其中W = Σw\_m为总权重，I(·)为指示函数，P\_joint(m)为任务点m处的联合探测概率，P\_th为探测阈值。

联合探测概率采用OR融合规则：
\begin{equation}
    P_{\text{joint}}(m) = 1 - \prod_{j=1}^{J} (1 - P_{\text{detect}}(j, m))
\end{equation}

探测概率模型：
\begin{equation}
    P_{\text{detect}}(j, m) = P_0 \cdot \exp(-\beta \cdot d_{jm})
\end{equation}
其中P\_0为最大探测概率，beta为衰减系数，d\_jm为雷达j与任务点m的欧氏距离。

\textbf{目标函数2 - 压制效能：}
\begin{equation}
    J_{\text{min}} = \min_{m \in T} \sum_{j=1}^{J} J_j(m)
\end{equation}
其中J\_j(m)为第j个干扰源在任务点m处的干扰功率密度。

干扰功率密度模型：
\begin{equation}
    J_j(m) = \frac{P_j}{4\pi d_{jm}^{\alpha_j}}
\end{equation}
其中P\_j为干扰源j的发射功率，lpha\_j为路径损耗指数(A2G: 2.0, G2G: 4.0)。

为将最大化问题转化为最小化问题，令f1 = 1 - ECR，f2 = 1/J\_min。优化目标为最小化(f1, f2)。

\subsection{区域分解算法}

采用Hertel-Mehlhorn算法进行区域凸分解。该算法的时间复杂度为O(n log n)，其中n为多边形顶点数。

算法步骤：
\begin{enumerate}
    \item \textbf{连通性处理}：将输入的多边形分解为连通分量
    \item \textbf{空洞消除}：使用两条线段切割方法消除所有空洞
    \item \textbf{三角剖分}：对无空洞多边形进行三角剖分
    \item \textbf{三角形合并}：合并相邻三角形形成凸多边形
\end{enumerate}

每个凸多边形被分配唯一的二进制编码，便于MOPSO优化。

\subsection{坐标变换}

使用垂直交点法进行坐标变换。给定凸多边形P和归一化坐标(ĥx, ĥy) in [0,1]²，变换步骤：

\begin{algorithm}
\caption{垂直交点坐标变换}
\begin{algorithmic}
\STATE 输入：凸多边形P，归一化坐标(ĥx, ĥy)
\STATE 在P的边界上找到与垂直线x = ĥx·width的交点
\STATE 在P的边界上找到与水平线y = ĥy·height的交点
\STATE 选择距离(ĥx, ĥy)最近的交点作为投影点Q
\STATE 将(ĥx, ĥy)映射到Q附近的小区域
\STATE 输出：物理坐标(x, y)
\end{algorithmic}
\end{algorithm}

该变换确保归一化坐标空间均匀映射到物理空间，有效消除边界效应。

\subsection{MOPSO-DT优化}

粒子编码：
\begin{equation}
    \Phi_j = [\hat{x}_j, \hat{y}_j, b_j^1, b_j^2, ..., b_j^{N_{\text{bin}}}]
\end{equation}
其中(ĥx\_j, ĥy\_j) in [0,1]²为归一化坐标，b\_j为凸多边形索引的二进制编码。

速度更新：
\begin{equation}
    v_i^{t+1} = w \cdot v_i^t + c_1 \cdot r_1 \cdot (p_i^{\text{best}} - x_i^t) + c_2 \cdot r_2 \cdot (g^{\text{best}} - x_i^t)
\end{equation}

其中w为惯性权重，c\_1和c\_2为学习因子，r\_1和r\_2为均匀随机数。

使用Pareto支配关系确定个体最优和全局最优。对于多目标问题，全局最优从外部档案中随机选择。

拥挤距离计算：
\begin{equation}
    d_i = \frac{f_1(x_{i+1}) - f_1(x_{i-1})}{f_1^{\max} - f_1^{\min}} + \frac{f_2(x_{i+1}) - f_2(x_{i-1})}{f_2^{\max} - f_2^{\min}}
\end{equation}

归一化目标函数：
\begin{equation}
    f_1^{\text{norm}} = 1 - \text{ECR}, \quad f_2^{\text{norm}} = \frac{J_{\text{min}}}{J_{\text{min}} + J_{\text{max}}^{\text{ref}}}
\end{equation}